\noindent{\heiti 摘要：}卫星海洋遥感大数据为自主水下航行器提供了大尺度先验环境场，使机动平台能够面向高动态海洋现象开展靶向观测。然而水下作业期间缺乏实时通信且机载电池与任务时间受限，面对滞后数日且缺失小尺度结构的卫星底图，航行器必须在开启高功耗芯片深入规划航线、维持低功耗处理器后台规划继续巡航采样以及关停推进器原位复测之间合理分配时间与能量。现有固定分配策略依赖同型平台的历史任务调参，而经典在线单步元推理直接用规划器在搜索样本上的内部得分增量评估计算收益，会因优化器诅咒高估计算价值并陷入过度规划。为此，本文提出以卫星大数据标定的留出想象海洋作为评判计算价值的独立标尺，构建出发前全任务预演与在线留出估值方法。该方法利用八年东海卫星叶绿素数据标定滞后底图误差，围绕3天前旧图生成48个随现场采样集合同化更新的想象海洋，并将其固定划分为32个搜索用海洋与16个留出海洋。下水前由母船或岸基计算机以4个留出海洋为假想真值，按机载算力与电池功耗完整模拟五种用时策略的24小时任务并选出时间平均遗憾最低者，或在线按周期在留出海洋上评估计算与采样收益。基于2018至2026年东海3197天数据构建的422个赤潮任务窗口、四种机载硬件及两种真值模型的实验表明，留出海洋评判在64至200,000次推演下与真实航线价值偏差不超过0.005。在148个测试任务上，出发前预演在主真值设置下的时间平均遗憾为0.265至0.276，无需历史任务即达到调好的固定策略与按历史任务学习的选择器水平，并比经典单步元推理降低0.011至0.022，同时在线引入留出估值可将单步元推理开启高功耗芯片的比例从80\%至97\%压低至30\%至44\%并达到预演水平。上述结果证实，将历史卫星遥感数据转化为独立于搜索过程的留出想象海洋可有效消除计算自评偏差，使航行器在无历史巡测记录条件下实现机载算力与观测行动的自适应权衡。

\vspace{4pt}
\noindent{\heiti 关键词：}海洋大数据；自主水下航行器；元推理；集合数据同化；赤潮巡测；优化器诅咒

\vspace{10pt}
\begin{center}
  {\zihao{4}\bfseries Balancing Computation and Observation for Autonomous Underwater Vehicles via Data-Calibrated Held-Out Imagined Oceans and Pre-Launch Rehearsal\par}
  \vspace{4pt}
  {\zihao{5} XU Chendong\par}
  {\zihao{-5}\itshape (Ocean College, Zhejiang University, Zhoushan 316021, China)\par}
\end{center}
\vspace{4pt}

\noindent{\bfseries Abstract:} Satellite ocean remote sensing big data provides large-scale prior fields for autonomous underwater vehicles, enabling targeted mobile observation of dynamic marine phenomena. Because underwater operations lack real-time communication and are constrained by onboard battery capacity and mission duration, a vehicle carrying a multi-day-old, spatially smoothed satellite map must continuously trade off activating a high-power chip for deep route planning against cruising to collect new samples under low-power background planning or stopping propulsion to resample locally. Existing fixed time-allocation rules require tuning on historical missions from the same hardware platform, whereas classic online one-step metareasoning evaluates computation using the planner's own score improvement on the searched ensemble members, which is inflated by the optimizer's curse and leads to severe over-planning. This paper proposes using uncertainty-calibrated held-out imagined oceans derived from satellite big data as an independent yardstick to judge the value of computation, instantiated both as pre-launch whole-mission rehearsal and as online held-out valuation. Using eight years of East China Sea satellite chlorophyll-a data to calibrate stale-map error, the vehicle generates 48 imagined oceans around a 3-day-old map, updates them with in-situ fluorometer measurements via a localised ensemble adjustment Kalman filter, and partitions them into 32 searched oceans and 16 held-out oceans. Before launch, a support-ship or shore computer treats 4 held-out oceans as imagined truths and simulates complete 24-hour missions for five time-allocation strategies under the vehicle's actual processor speed and battery budget to select the strategy with the lowest mean regret, while the online variant evaluates single-step computation and sampling gains on the 16 held-out oceans. Experiments across 422 bloom windows extracted from 3,197 daily maps (2018--2026), four hardware configurations, and two truth models show that held-out evaluation remains within 0.005 of true route value from 64 to 200,000 rollouts. Across 148 test missions under the main truth model, pre-launch rehearsal achieves time-averaged regrets of 0.265--0.276, matching fixed strategies tuned on historical missions and a learned strategy selector without requiring prior missions, while outperforming classic one-step metareasoning by 0.011--0.022 and heuristic survey rules. Replacing self-valuation with held-out valuation inside online one-step metareasoning similarly reduces high-power chip activation from 80\%--97\% to 30\%--44\% of intervals and matches rehearsal performance. These results demonstrate that decoupling route search from value judgement via data-calibrated held-out imagined oceans eliminates planner self-assessment bias and adapts time allocation to onboard hardware without historical mission tuning.

\vspace{4pt}
\noindent{\bfseries Keywords:} ocean big data; autonomous underwater vehicle; metareasoning; ensemble data assimilation; harmful algal bloom survey; optimizer's curse
